% Options for packages loaded elsewhere
\PassOptionsToPackage{unicode}{hyperref}
\PassOptionsToPackage{hyphens}{url}
\PassOptionsToPackage{space}{xeCJK}
%
\documentclass[
]{article}
\usepackage{amsmath,amssymb}
% \usepackage{lmodern}
\usepackage{iftex}
\ifPDFTeX
  \usepackage[T1]{fontenc}
  \usepackage[utf8]{inputenc}
  \usepackage{textcomp} % provide euro and other symbols
\else % if luatex or xetex
  \usepackage{unicode-math}
  \defaultfontfeatures{Scale=MatchLowercase}
  \defaultfontfeatures[\rmfamily]{Ligatures=TeX,Scale=1}
  \setmainfont[]{DejaVu Serif}
  \setmathfont[]{Asana Math}
  \ifXeTeX
    \usepackage{xeCJK}
    \setCJKmainfont[]{Noto Serif CJK TC}
  \fi
  \ifLuaTeX
    \usepackage[]{luatexja-fontspec}
    \setmainjfont[]{Noto Serif CJK TC}
  \fi
\fi
% Use upquote if available, for straight quotes in verbatim environments
\IfFileExists{upquote.sty}{\usepackage{upquote}}{}
\IfFileExists{microtype.sty}{% use microtype if available
  \usepackage[]{microtype}
  \UseMicrotypeSet[protrusion]{basicmath} % disable protrusion for tt fonts
}{}
\makeatletter
\@ifundefined{KOMAClassName}{% if non-KOMA class
  \IfFileExists{parskip.sty}{%
    \usepackage{parskip}
  }{% else
    \setlength{\parindent}{0pt}
    \setlength{\parskip}{6pt plus 2pt minus 1pt}}
}{% if KOMA class
  \KOMAoptions{parskip=half}}
\makeatother
\usepackage{xcolor}
\IfFileExists{xurl.sty}{\usepackage{xurl}}{} % add URL line breaks if available
\IfFileExists{bookmark.sty}{\usepackage{bookmark}}{\usepackage{hyperref}}
\hypersetup{
  hidelinks,
  pdfcreator={LaTeX via pandoc}}
\urlstyle{same} % disable monospaced font for URLs
\setlength{\emergencystretch}{3em} % prevent overfull lines
\providecommand{\tightlist}{%
  \setlength{\itemsep}{0pt}\setlength{\parskip}{0pt}}
\setcounter{secnumdepth}{-\maxdimen} % remove section numbering
\usepackage[margin=1in]{geometry}
\usepackage{microtype}
\usepackage{hyperref}
\hypersetup{colorlinks=true, linkcolor=blue, urlcolor=blue}

\usepackage{fontspec}
\usepackage{xeCJK}

\setmainfont{DejaVu Serif}
\setCJKmainfont{Noto Serif CJK TC}[BoldFont = Noto Serif CJK TC]

\usepackage{amsmath, amssymb, amsthm}
\ifLuaTeX
  \usepackage{selnolig}  % disable illegal ligatures
\fi

\author{}
\date{}

\begin{document}

HORS（Hash to Obtain Random
Subsets）是一種少數次簽章方案，它利用哈希函數和隨機子集的概念，實現了簽章的快速生成與驗證。同時，通過適當的參數選擇，HORS
可以在安全性和效率之間取得良好的平衡。

\hypertarget{ux7cfbux7d71ux4ecbux7d39}{%
\section{系統介紹}\label{ux7cfbux7d71ux4ecbux7d39}}

我們首先將要簽章的訊息分為區段： \[
d = d_{1},d_{2},\dots,d_{end}
\] 每段都有 log2(t) 位元（其中 t 為系統參數，通常是 2 的某次方）

接著將每個區段解釋為一個二進位整數，叫做索引值。 \[
\mathrm{indices} = \{ \mathrm{index}_{1},\mathrm{index}_{2},\dots, \mathrm{index}_{end}\}
\] \[
\mathrm{index}_{i} = \mathrm{int}(d_{i},2)
\] （意思是將區塊 d\_i 當作二進位整數，並輸出為一個整數 index\_i）

生成 t 個私鑰，每個私鑰都是 n 位元的二進位字串，而他們的雜湊值就是公鑰。
\[
\mathrm{SK}=\{\mathrm{sk}_{0},\mathrm{sk}_{1},\dots,\mathrm{sk}_{t-1}\}
\]\[
\mathrm{PK}=\{\mathrm{pk}_{0},\mathrm{pk}_{1},\dots,\mathrm{pk}_{t-1}\}
\] \[
\mathrm{pk}_{i}= \mathcal{H}(\mathrm{sk}_{i})
\] 最後，就可以設計簽章為 \[
\sigma 
= \{\sigma_{1},\dots,\sigma_{end}\}
\] \[
\sigma_{i} = \mathrm{sk}_{\mathrm{index}_{i}} = \mathrm{SK}[\mathrm{index}_{i}]
\] 而驗章時就對每個區段檢查 \[
\mathcal{H}(\sigma_{i}) = \mathrm{PK}[\mathrm{index}_{i}]
\]

\hypertarget{ux5b8cux6574ux7684ux7cfbux7d71}{%
\section{完整的系統}\label{ux5b8cux6574ux7684ux7cfbux7d71}}

\hypertarget{ux53c3ux6578ux8207ux9470ux5319ux751fux6210}{%
\subsection{參數與鑰匙生成}\label{ux53c3ux6578ux8207ux9470ux5319ux751fux6210}}

\hypertarget{ux53c3ux6578}{%
\subsubsection{參數}\label{ux53c3ux6578}}

參數 t ：我們會把訊息分割為 log2(t) 的長度 \[
d = \underbrace{ d_{1} }_{ \log_{2}t },\underbrace{ d_{2} }_{ \log_{2}t},\dots,\underbrace{ d_{end} }_{ \log_{2}t}
\] 我們設定所使用的雜湊函數是使用 SHA-256
，也設定這個訊息是雜湊函數後的結果（因此固定長度）

\hypertarget{ux516cux79c1ux9470}{%
\subsubsection{公私鑰}\label{ux516cux79c1ux9470}}

私鑰 \[
\mathrm{SK}=\{\mathrm{sk}_{0},\mathrm{sk}_{1},\dots,\mathrm{sk}_{t-1}\}
\] 其中每個 sk 都是長度為 n = 256 的二進位字串。 \[
\mathrm{PK}=\{\mathrm{pk}_{0},\mathrm{pk}_{1},\dots,\mathrm{pk}_{t-1}\}
\] \[
\mathrm{pk}_{i}= \mathcal{H}(\mathrm{sk}_{i})
\]

因為我們選用 SHA-256 所以每個 pk 都是 n = 256 長度的二進位字串。
公佈公鑰給驗章者

\hypertarget{ux7c3dux7ae0}{%
\subsection{簽章}\label{ux7c3dux7ae0}}

首先先進行訊息分段： \[
d = \underbrace{ d_{1} }_{ \log_{2}t },\underbrace{ d_{2} }_{ \log_{2}t},\dots,\underbrace{ d_{end} }_{ \log_{2}t}
\] 計算每個區塊的索引 \[
\mathrm{index}_{i} = \mathrm{int}(d_{i},2)
\] \[
\mathrm{indices} = \{ \mathrm{index}_{1},\mathrm{index}_{2},\dots, \mathrm{index}_{end}\}
\]

並找出每段對應的私鑰 \[
\sigma 
= \{\sigma_{1},\dots,\sigma_{end}\}
\] \[
\sigma_{i} = \mathrm{sk}_{\mathrm{index}_{i}} = \mathrm{SK}[\mathrm{index}_{i}]
\]

發送訊息跟簽章給驗章者 \#\# 驗章

首先先進行訊息分段： \[
d = \underbrace{ d_{1} }_{ \log_{2}t },\underbrace{ d_{2} }_{ \log_{2}t},\dots,\underbrace{ d_{end} }_{ \log_{2}t}
\] 計算每個區塊的索引 \[
\mathrm{index}_{i} = \mathrm{int}(d_{i},2)
\] \[
\mathrm{indices} = \{ \mathrm{index}_{1},\mathrm{index}_{2},\dots, \mathrm{index}_{end}\}
\] 驗證所有索引的 hash 值 \[
\mathcal{H}(\sigma_{i}) 
\overset{?}{=} \mathrm{pk}_{\mathrm{index}_{i}} = \mathrm{PK}[\mathrm{index_{i}}]
\] 如果所有的等號都成立，則驗章成功。

\hypertarget{ux5be6ux4f5cux7cfbux7d71}{%
\section{實作系統}\label{ux5be6ux4f5cux7cfbux7d71}}

首先定義好雜湊函數：

\begin{verbatim}
import hashlib
import random

# 定義哈希函數 (SHA-256)，輸入和輸出均為位元串
def hash_func(bitstring):
    # 將位元串轉換為位元組
    byte_data = int(bitstring, 2).to_bytes((len(bitstring) + 7) // 8, byteorder='big')
    # 計算 SHA-256 哈希值
    sha256_hash = hashlib.sha256(byte_data).digest()
    # 將哈希結果轉換為二進位格式，並返回 256 位的二進位字串
    return ''.join(f'{byte:08b}' for byte in sha256_hash)
\end{verbatim}

\hypertarget{ux53c3ux6578ux8207ux9470ux5319ux751fux6210-1}{%
\subsection{參數與鑰匙生成}\label{ux53c3ux6578ux8207ux9470ux5319ux751fux6210-1}}

\begin{verbatim}
# 參數選擇
logt = 8
t = 2**8  # 公鑰中哈希值的總數量
k = 32     # 簽章中選擇的哈希值數量
n = 256    # 哈希輸出的位元長度（256 位元）

# 生成 n 位元的隨機位串
def generate_random_bitstring(n):
    return ''.join(str(random.randint(0, 1)) for _ in range(n))

# 密鑰生成
def keygen():
    SK = [generate_random_bitstring(n) for _ in range(t)]
    PK = [hash_func(sk) for sk in SK]
    return SK, PK

SK,PK = keygen()
print("SK[0]:", SK[0][:64], "...")
print("PK[0]:", PK[0][:64], "...")


# Outputs:
# SK[0]: 1010010010111001011010101110111100001101010000100011011001101000 ...
# PK[0]: 0011010010110000001001011000100011010010010111110000100001010110 ...
\end{verbatim}

\hypertarget{ux7c3dux7ae0-1}{%
\subsection{簽章}\label{ux7c3dux7ae0-1}}

\begin{verbatim}
# 將訊息轉換為他的 hash 值

def message_to_binary(message):
    return ''.join(format(ord(c), '08b') for c in message)

message = "Hi! This is Cesare!"
message_bits = message_to_binary(message)
message_hash = hash_func(message_bits)

print(message_hash)

# Outputs:
# 0010111001001011011001000011001001111100000000111010011110001010010011110011001000011001111010100001110001101110101101001001110001100111101101011001000101010011011101111100111101011001111110100100000010001101010100111000001010010101010011011101101101001111
\end{verbatim}

\begin{verbatim}
# 將訊息的哈希值分區段，並將其解釋為整數

def message_to_indices(message_hash):
    indices = []
    for ptr in range(0,n,8):
        indices.append(int(message_hash[ptr:ptr+8],2))
    return indices

indices = message_to_indices(message_hash)
print(indices)

# Outputs:
# [46, 75, 100, 50, 124, 3, 167, 138, 79, 50, 25, 234, 28, 110, 180, 156, 103, 181, 145, 83, 119, 207, 89, 250, 64, 141, 83, 130, 149, 77, 219, 79]
\end{verbatim}

\begin{verbatim}
# 生成簽章

signature = ""
for index in indices:
    signature += SK[index]

print(signature[:2*n])

# Outputs:
# 11110111001110000...10001101
\end{verbatim}

\hypertarget{ux9a57ux7ae0}{%
\subsection{驗章}\label{ux9a57ux7ae0}}

\begin{verbatim}
# 驗證簽章
# 驗證者有 message, signature


message_hash = hash_func(message_bits)

indices = []
for ptr in range(0,n,8):
    indices.append(int(message_hash[ptr:ptr+8],2))


validity = True
for i in range(n//logt):
    index = indices[i]
    signature_chunk = signature[i*n:(i+1)*n]

    if not hash_func(signature_chunk)==PK[index]:
        print("The signature is not valid!")
        validity = False
        break

print(validity)

# Outputs:
# True
\end{verbatim}

\hypertarget{ux6e2cux8a66ux7ac4ux6539ux6d88ux606fux7684ux653bux64ca}{%
\section{測試竄改消息的攻擊}\label{ux6e2cux8a66ux7ac4ux6539ux6d88ux606fux7684ux653bux64ca}}

\begin{verbatim}
# 篡改消息
tampered_message = "Hi! This is Caesar!"  # 修改了消息內容
tampered_message_bits = message_to_binary(tampered_message)
tampered_message_hash = hash_func(tampered_message_bits)
tampered_indices = message_to_indices(tampered_message_hash)

# 驗證篡改的消息
validity_tampered = True
for i in range(256 // logt):
    index = tampered_indices[i]
    signature_chunk = signature[i * n:(i + 1) * n]
    if hash_func(signature_chunk) != PK[index]:
        print("The tampered signature is not valid!")
        validity_tampered = False
        break

print("篡改消息的驗證結果：", validity_tampered)

# Outputs:
# The tampered signature is not valid!
# 篡改消息的驗證結果： False
\end{verbatim}

ref SRIVASTAVA, Vikas; BAKSI, Anubhab; DEBNATH, Sumit Kumar. An overview
of hash based signatures.~\emph{Cryptology ePrint Archive}, 2023.

\end{document}
